\chapter{Quantum Parallel Tempering (\code{sqapt})}
\label{ch:sqapt}

\code{sqapt} combines the two strongest ideas of the previous chapters:
the Suzuki--Trotter quantum simulation of Chapter~\ref{ch:sqa} and the
replica-exchange machinery of Section~\ref{sec:classical-pt}. It is the
\textbf{recommended default} for rugged, multi-modal problems: several
complete SQA systems run simultaneously at different points of a
$(\beta,\Gamma)$ ladder, and adjacent systems periodically swap
configurations so that solutions discovered under strong fluctuations flow
toward the strongly-frozen rungs where they are polished.

\section{The $(\beta,\Gamma)$ ladder}

Each of the $R$ replicas $r=1,\dots,R$ is a full Trotter system
($n\times M$ spins) equilibrating at its own pair $(\beta_r,\Gamma_r)$,
hence its own action $S_r$ (Eq.~\eqref{eq:action}) with its own
$\Jp^{(r)}=\Jp(\beta_r,\Gamma_r,M)$. The ladder is arranged so that mobility
decreases along it:

\begin{center}
\begin{tikzpicture}[scale=1.0,
  rung/.style={rectangle,rounded corners=1.5pt,draw=qnavy!60,fill=qnavy!7,
               minimum width=3.35cm,minimum height=0.72cm,font=\small},
  sw/.style={{Stealth}-{Stealth},thick,qteal}]
\node[rung,fill=qamber!12,draw=qamber!70] (r1) at (0,0)
  {$r{=}1$: hot, quantum};
\node[rung] (r2) at (4.1,0) {$r{=}2$};
\node[font=\small\color{qgray}] (dots) at (7.0,0) {$\cdots$};
\node[rung,fill=qnavy!16] (rR) at (9.9,0) {$r{=}R$: cold, classical};
\draw[sw] (r1) -- node[above,font=\scriptsize\sffamily\color{qteal}]{swap} (r2);
\draw[sw] (r2) -- (dots);
\draw[sw] (dots) -- node[above,font=\scriptsize\sffamily\color{qteal}]{swap} (rR);
\node[below=1.5mm of r1,font=\scriptsize\sffamily\color{qgray}]
  {small $\beta$, large $\Gamma$};
\node[below=1.5mm of rR,font=\scriptsize\sffamily\color{qgray}]
  {large $\beta$, small $\Gamma$};
\end{tikzpicture}
\end{center}

\code{auto\_ladder\_sqa\_tuned(problem, replicas=R, mode=...)} builds the
ladder from the problem's probed energy scale:
\begin{itemize}
\item $\beta$ rungs are spaced \emph{geometrically} from
$\beta_{\min}$ (hot) to $\beta_{\max}$ (cold), so adjacent-rung energy
distributions overlap roughly uniformly along the ladder;
\item $\Gamma$ rungs are spaced \emph{inverse-geometrically} --- high
$\Gamma$ where $\beta$ is small, low $\Gamma$ where $\beta$ is large --- so
every rung sits at a physically coherent point (thermal and quantum mobility
rise and fall together).
\end{itemize}

\section{The quantum swap rule, derived}
\label{sec:quantum-swap}

The joint stationary measure is
$\Pi(\{x_r\})\propto\prod_r\mathrm{e}^{-S_r(x_r)}$, where $x_r$ denotes the
entire Trotter configuration of replica $r$ and
\begin{equation}
S_r(x)=\frac{\beta_r}{M}\sum_{t}\Ecl(s^t)\;-\;\Jp^{(r)}\sum_{t,i}s^t_i s^{t+1}_i
\;\equiv\;\frac{\beta_r}{M}\,C(x)\;-\;\Jp^{(r)}\,B(x),
\label{eq:action-cb}
\end{equation}
with $C(x)$ the summed classical energy over slices and $B(x)$ the total
imaginary-time \emph{bond sum}. (The quantity $B$ returns as the central
observable of Chapter~\ref{ch:optimal}.)

A swap proposes exchanging the full configurations of rungs $r$ and $r{+}1$.
Detailed balance for $\Pi$ gives the acceptance

\begin{keyeq}[SQAPT swap acceptance]
\begin{equation}
A_{\text{swap}}=\min\Bigl(1,\;
\exp\bigl[-\bigl(S_r(x_{r+1})+S_{r+1}(x_r)-S_r(x_r)-S_{r+1}(x_{r+1})\bigr)\bigr]\Bigr).
\label{eq:sqapt-swap}
\end{equation}
\end{keyeq}

This is precisely what the C++ core computes: it evaluates the four
cross-actions $S_r(x_r)$, $S_{r+1}(x_{r+1})$, $S_r(x_{r+1})$,
$S_{r+1}(x_r)$ from Eq.~\eqref{eq:action-cb} and applies
Eq.~\eqref{eq:sqapt-swap}. Expanding via Eq.~\eqref{eq:action-cb}, the
exponent splits into a thermal part and a quantum part:
\begin{equation}
\Delta_{\text{swap}}
=\frac{\beta_{r}-\beta_{r+1}}{M}\bigl(C(x_{r+1})-C(x_{r})\bigr)
\;-\;\bigl(\Jp^{(r)}-\Jp^{(r+1)}\bigr)\bigl(B(x_{r+1})-B(x_{r})\bigr).
\label{eq:swap-split}
\end{equation}
The first term is the classical PT criterion Eq.~\eqref{eq:pt-swap} applied
to slice-summed energies; the second exchanges configurations according to
their quantum ``kinetic'' content. When all replicas \emph{share} one
$\Jp$ --- as in the optimal-schedule mode of
Section~\ref{sec:sqapt-optimal} --- the second term vanishes identically and
the swap reduces to a purely classical criterion.

\section{The algorithm}

\code{SQAParallelTemperingAnnealer.run()} iterates \code{steps} epochs; in
each epoch every replica performs its local SQA updates
(\code{sweeps\_per\_step} single-spin sweeps, \code{worldline\_sweeps}
worldline sweeps, optionally \code{cluster\_sweeps} SW sweeps --- all with
the rules of Section~\ref{sec:sqa-moves} at that replica's own
$(\beta_r,\Gamma_r,\Jp^{(r)})$), and every \code{swap\_interval} epochs a
sweep of adjacent-pair swaps is attempted with
Eq.~\eqref{eq:sqapt-swap}. Replicas are OpenMP-parallel; the per-epoch swap
acceptance fraction is recorded in \code{swap\_acceptance\_trace}.

Unlike annealing, the ladder is \emph{static}: no schedule ramps during the
run. Optimisation progress comes from configurations drifting down the
ladder. The best classical projection (best single slice, any replica, any
epoch) is tracked exactly as in SQA.

\section{Usage}

\subsection{One call}

\begin{pycode}
from qanneal import solve

result = solve(
    ising,
    method="sqapt",
    replicas=8,          # ladder rungs
    reads=8,
    trotter_slices=24,
    sweeps_per_beta=50,  # local sweeps per epoch
    worldline_sweeps=4,
    cluster_sweeps=1,    # SQAPT+SW
    pt_steps=70,         # epochs
    swap_interval=1,
    seed=3,
)
\end{pycode}

If no schedule/ladder is supplied, \code{solve()} calls
\code{auto\_ladder\_sqa\_tuned} internally. Explicit ladders:

\begin{pycode}
from qanneal import auto_ladder_sqa_tuned, SQAParallelTemperingAnnealer

ladder = auto_ladder_sqa_tuned(ising, replicas=8, mode="balanced")
ann = SQAParallelTemperingAnnealer(
    ising, ladder.betas, ladder.gammas, trotter_slices=24)
res = ann.run(sweeps_per_step=50, worldline_sweeps=4,
              steps=80, swap_interval=1, cluster_sweeps=1)
print(res.best_energy)
print(res.swap_acceptance_trace[-5:])   # aim for ~0.2-0.5
\end{pycode}

You may also hand explicit ladders to \code{solve()} via
\code{pt\_betas=[...]}, \code{pt\_gammas=[...]} (they must pair up).

\subsection{Reading the diagnostics}

\begin{itemize}
\item \code{swap\_acceptance\_trace} --- mean accepted-swap fraction per
epoch. Healthy: $0.2$--$0.5$. Too low $\Rightarrow$ add rungs or narrow the
ladder; too high $\Rightarrow$ rungs redundant.
\item \code{final\_energies} --- final projected energy of each rung; should
be ordered roughly hot-to-cold. A cold rung stuck far above the best energy
signals insufficient epochs or a mobility bottleneck mid-ladder.
\item \code{average\_energy\_trace} --- ladder-mean energy per epoch;
plateaus indicate equilibration.
\end{itemize}

\section{Choosing \code{sqapt}}

\begin{itemize}
\item \textbf{Use it} for frustrated, glassy, multi-modal problems ---
anything where single-trajectory annealing shows large read-to-read
variance.
\item Ladder size \code{replicas}$\,=8$ is a solid default; scale up with
problem hardness, not with $n$.
\item \code{cluster\_sweeps=1} (SQAPT+SW) is the strongest general-purpose
configuration in the library.
\item Prefer \code{sqa} when the problem is easy enough that restarts
suffice --- SQAPT's ladder overhead ($R$ full Trotter systems) is then wasted.
\end{itemize}

The combination of \code{sqapt} with the adaptive $\Jp$ schedule --- a
shared, susceptibility-paced Trotter coupling ramp across the whole ladder ---
is covered in Section~\ref{sec:sqapt-optimal}.
